%% Sections shared by every manual, English edition.
%% Pulled in with \input{common-en}. The only thing that differs between
%% applets is \appletfolder, which is also the folder name and the path on the site.

\section{Opening it}

There is nothing to install. The whole application is one \code{index.html}
file: it downloads no libraries, calls no server, and stores nothing outside
your browser.

\subsection{The three ways}

\begin{description}
\item[Double-click the file.] The quickest route. Find
  \code{\appletfolder/index.html} and open it. The browser shows it straight off the
  disk, with an address beginning \code{file://}. It works offline from the
  first moment.
\item[From the web.] If someone has published it, the address ends in
  \code{/\appletfolder/}. Once the page has loaded you can disconnect: it needs the
  network for nothing else.
\item[Served from a folder.] For a class, putting the folder on any static
  server is enough. With Python installed:
  \begin{quote}\code{python3 -m http.server 8000}\end{quote}
  then \code{http://localhost:8000/\appletfolder/} in the browser.
\end{description}

\begin{amnote}
Saving the page with \key{Ctrl}\,+\,\key{S} (or \key{Cmd}\,+\,\key{S} on a
Mac) leaves a copy of your own that still works offline and that can travel on a
memory stick or go out by email.
\end{amnote}

\subsection{Language and theme}

Two pairs of buttons sit at the top right.

\begin{ctrltable}{Control & What it does}
ES / EN & Switches the whole interface between Spanish and English on the spot.
          Nothing is recomputed, so you can switch mid-explanation without
          losing the state you had. \\
Dark / Light & Switches the theme. Dark is the original design; light exists
          for projectors and for printing. \\
\end{ctrltable}

Both choices are remembered in the browser, so next time it opens the way you
left it.

\subsection{Which browser}

Any reasonably recent one: Chrome, Edge, Firefox or Safari, on a computer,
tablet or phone. On a narrow screen the panels stack one above the other
instead of sitting side by side. No account, no permissions, no extensions.
